\chapter{Análise do Script MATLAB: Hasse}

\section{Visão Geral}

\textbf{What is the main purpose of this file?}
The main purpose of the \texttt{hasse.m} script is to compute the spatial coordinates (x and y) and plot the edges to visually render a Hasse diagram. A Hasse diagram is a mathematical visualization used to represent a finite partially ordered set.

\textbf{What type of analysis or calculation does it perform?}
The code performs matrix manipulations and path-length calculations. By multiplying the adjacency matrix by itself, it calculates the levels (y-coordinates) of each node based on the longest paths from the minimal elements. It also calculates horizontal placement (x-coordinates) to evenly space elements that share the same y-coordinate layer.

\textbf{What is the mathematical/theoretical context?}
The context relies heavily on \textit{Order Theory} and \textit{Grafo Theory}. It deals with \textit{binary relations}, specifically partial orders, which must be reflexive, transitive, and antisymmetric. The script translates the properties of Directed Acyclic Grafos (DAGs) into a structured visualization known as a Hasse diagram, using an adjacency matrix representation of the relation.

\vspace{0.5cm}
\hrule
\vspace{0.5cm}

\section{Análise Passo a Passo do Código}

\subsection*{1. Título da Secção/Função: Entrada Parsing and Initialization}
\textbf{2. Explanation:} This block manages the input arguments provided to the function. It establishes the elements D of the relation. If the user passes a 1-by-1 matrix (a single scalar n), it automatically interprets this as a request to build a divisibility relation up to n. Otherwise, it defines D based on the size of the adjacency matrix R or reads D directly from the optional arguments (\texttt{varargin}).

\textbf{3. Code Snippet:}
\begin{lstlisting}
if isequal(size(R),[1 1])
div=@(n) find(mod(n,1:n)==0);
n=R;
D=div(n), [i,j]=ndgrid(1:numel(D)); R=(mod(D(j),D(i))==0);
elseif nargin==1,
D=1:size(R,1);
else
D=varargin{1};
end
\end{lstlisting}

\vspace{0.5cm}

\subsection*{1. Título da Secção/Função: Mathematical Validation of the Partial Order}
\textbf{2. Explanation:} Before computing the layout, the script validates that the provided matrix R strictly represents a partial order. It checks three conditions using matrix logic:

Transitivity (\texttt{logical(R*R)<=R})

Reflexivity (\texttt{diag(R)}) meaning all diagonal elements are 1.

Antisymmetry (\texttt{(R\&R')<=eye(size(R))}).
If any of these conditions fail, it throws a logical error halting execution.

\textbf{3. Code Snippet:}
\begin{lstlisting}
nD=numel(D);
if ~isequal(size(R),[nD nD]), error('R must be nD-by-nD, with nD=numel(D)'), end
if ~all(all(logical(R*R)<=R)) | ~all(diag(R)) | ~all(all((R&R')<=eye(size(R))))
error('R is not a partial order, i.e. a reflexive, transitive and antisymmetric relation')
end
\end{lstlisting}

\vspace{0.5cm}

\subsection*{1. Título da Secção/Função: Y-Coordinate (Level) Calculation}
\textbf{2. Explanation:} This section calculates the vertical level (y-coordinate) for each node. It isolates edges of exactly length 1 (\texttt{R1}), ignoring transitivity leaps. It increments the y value of a target node by tracking paths. By repeatedly multiplying \texttt{Rn} by \texttt{R1}, it finds paths of length n. The loop runs as long as paths of length n exist, effectively pushing a node up by 1 level for every step in the longest chain leading to it.

\textbf{3. Code Snippet:}
\begin{lstlisting}
y=ones(nD,1);
R1=RR==2;
[d,c]=find(R1);
y(c)=y(c)+1;
Rn=R1R1;
while nnz(Rn)>0
[i,j]=find(Rn);
y(j)=y(j)+1;
Rn=Rn*R1;
end
\end{lstlisting}

\vspace{0.5cm}

\subsection*{1. Título da Secção/Função: X-Coordinate Calculation}
\textbf{2. Explanation:} Once the vertical levels are established, nodes on the same level must be spaced out horizontally to prevent overlap. It sorts the nodes by their y values, counts how many nodes are on each level using \texttt{sparse}, and then distributes them symmetrically along the x-axis. It calculates a cumulative width and centers the nodes using \texttt{w(ys)/2}. A complex vector v=x+iy is also constructed representing the 2D plane coordinate for each node.

\textbf{3. Code Snippet:}
\begin{lstlisting}
w=sparse(y,1,1);
[ys,sy]=sort(y);
cw=cumsum(w);
x(sy)=(1:nD)'-cw(ys)+w(ys)/2;
x=x(:);
v=x+1i*y;
\end{lstlisting}

\vspace{0.5cm}

\subsection*{1. Título da Secção/Função: Visualization and Edge Plotting}
\textbf{2. Explanation:} The final step is to plot the graph onto a figure. It calculates the directional vectors (\texttt{du}, \texttt{dv}) from node d to node c. It uses MATLAB's \texttt{quiver} function to draw the lines without arrowheads, which is the standard convention for Hasse diagrams (direction is implied bottom-to-top). Finally, it overlays the text labels corresponding to the element values D at their respective (x,y) coordinates and turns off the axis grid for a cleaner mathematical aesthetic.

\textbf{3. Code Snippet:}
\begin{lstlisting}
du=x(c)-x(d);
dv=y(c)-y(d);
if ~isempty(d)
hq=quiver(x(d),y(d),du,dv,0);
set(hq,'showarrowhead','off')
end
text(x,y,num2str(D(:))), axis([min(x) max(x) min(y)-eps max(y)+eps]), axis off
\end{lstlisting}

\vspace{0.5cm}
\hrule
\vspace{0.5cm}

\section{Saídas e Elementos Visuais Esperados}

When this code executes successfully, it does not print values to the console but instead generates a \textbf{2D Grafoical Plot}.

\begin{itemize}
\item \textbf{The Plot Canvas:} The visual output is a standard MATLAB figure window containing a network/graph structure. The background is clean because the axes and tick marks are explicitly hidden (\texttt{axis off}).
\item \textbf{The Nodes:} The elements of the set D (e.g., numbers representing divisors) are rendered as raw text values suspended in the 2D space.
\item \textbf{The Y-Axis (Vertical Logic):} While invisible, the vertical axis represents the "hierarchy" or chain length of the partial order. Minimal elements (e.g., the number 1 in a divisibility relation) are at the very bottom. Elements that are ordered "higher" than others appear physically higher up on the canvas.
\item \textbf{The X-Axis (Horizontal Logic):} Elements that are incomparable but occupy the same hierarchical depth in the poset are distributed horizontally with equal spacing to minimize clustering.
\item \textbf{The Edges:} Solid lines connect the nodes. In adherence to Hasse diagram conventions, there are no arrowheads, and redundant transitive lines are omitted (e.g., if A connects to B, and B connects to C, the line from A directly to C is left out).
\end{itemize}

% Would you like me to help you debug any other parts of your MATLAB code,
% or perhaps convert this document into a different format?